\section{Exact source equivalence for guarded jump replacements}
\label{sec:zero-jump-source}

This certificate concerns the register program. It does not identify a
source instruction merely by the name of a physical decoder state.
A separate physical binding and blank-tape initialization are necessary
for any transition-table claim.

\paragraph{The finite invariant.}
A register vector $r\in\mathbb N^9$ has zero pattern
\[
 \alpha(r)=\sum_{i=0}^{8}2^i[r_i>0]\in\{0,\ldots,511\}.
\]
Let $P$ be the pinned 3,812-site PA196 register program, using the ordinary
cyclic source execution convention. The finite certificate stores sets
$J_i\subseteq\{0,\ldots,511\}$, one for each source site. Starting from
$(25,1)$, it closes under the following abstract transitions:
\begin{itemize}
\item Increment register $r$: set bit $r$ and take the fixed successor.
\item Decrement register $r$ with bit $r$ zero: preserve the mask and take
      the zero successor.
\item Decrement register $r$ with bit $r$ one: take the positive successor,
      and allow either retaining or clearing bit $r$.
\item Jump: preserve the mask and take its fixed successor.
\item Halt: no successor is required.
\end{itemize}
Each concrete step has one of the listed abstract successors. In the
positive decrement case the two possibilities correspond to a previous
value greater than one or equal to one. Thus these are all-value rules,
not observations of a finite concrete run.

The independently reconstructed closure contains exactly 274,880
site-mask pairs. The verifier checks the initial pair and every closure
edge explicitly, including under optimized Python execution. It also
checks the SHA-256 of the unchanged source file. The source control
successors are reconstructed from the literal binary-tree paths; they
are not accepted from an optimizer's receipt.

Define
\[
 \mathcal I(i,r)\quad\Longleftrightarrow\quad\alpha(r)\in J_i.
\]
Induction on source steps proves that $\mathcal I$ holds throughout the
run beginning at site 25 with $r_0=2$ and every other register zero. In
fact the same statement holds for any positive initial value of $r_0$,
because its initial zero pattern is the same.

\paragraph{Source site 26.}
The original instruction at path $100001$ is a forward jump to source
site 34. The finite invariant proves
\[
 \mathcal I(26,r)\ \Longrightarrow\ r_3=r_6=r_7=0.
\]
Replace that one source leaf by a decrement of $r_6$, with zero successor
34 and positive successor 27. Whenever $\mathcal I$ holds, the new
instruction takes its zero branch and leaves the entire register vector
unchanged. Its next configuration is therefore exactly the next
configuration of the original jump. This is one source step on either
side, so there is no time rescaling at source level.

\paragraph{Simultaneous-replacement theorem.}
Let $S$ be any finite collection of sites of $P$, where the instruction at
each $i\in S$ is a forward jump with successor $z_i$. Choose a register
$r(i)$ such that bit $r(i)$ is zero in every mask in $J_i$. Replace the
instruction by a decrement of $r(i)$ whose zero successor is $z_i$; its
positive successor may be arbitrary. Call the resulting program $P_S$.
Then for every configuration satisfying $\mathcal I$,
\[
       \operatorname{step}_{P_S}(i,r)=\operatorname{step}_{P}(i,r).
\]
The same $\mathcal I$ is inductive for $P_S$. Hence, from any common
initial configuration satisfying $\mathcal I$, the two source executions
agree at every step and halt at exactly the same source step, if either
halts.

The proof is pointwise: at an unchanged site the instructions agree; at
an edited site the tested register is zero, so both instructions perform
the same jump without modifying a register. Closure of $\mathcal I$
under the old instruction gives closure under the new instruction.
Induction then proves equality of complete executions. In particular,
several independently justified edits do not require assuming that the
new program continues to follow a sampled old trace: the old inductive
invariant remains valid for all of them simultaneously.

For $r(i)=6$, the certificate verifies the simultaneous replacement at
\[
 S_6=\{3,9,26,3665,3671,3683,3689,3692,3766,3772\}.
\]
The scripts retain both the single-site program and this ten-site
program as complete path/instruction JSON data. The other eligible
register choices are listed in the receipt.

\paragraph{PA formalizability.}
The zero-pattern map, the finite membership test for $J_i$, and both
source step functions are primitive recursive. Each fixed finite
certificate check can be proved in PA. The abstract-step soundness
argument is elementary arithmetic; the unbounded execution comparison
then follows by ordinary induction on the number of source steps.
Consequently this construction supplies a direct route to a PA proof
of the source-halting equivalence. No formalized PA proof or Lean
kernel certificate is asserted by the Python computation itself.

\paragraph{Why this does not validate the tempting 195-state table.}
The attempted change of native row $27/1$ to $1R62$ was informally
identified with decrementing register seven. The complete source-site
check disproved that identification: the decoder's effect depends on
its tape history. At the right-child path of source site 26 it does not
implement the required zero decrement and successor. A subsequent
locally certified fold of native states 27 and 46 can preserve the
behavior of this wrongly changed parent table without repairing its
source binding. Neither the source theorem above nor that local fold
therefore establishes the requested consistency equivalence for that
195-state table.

The source certificate and a physical runtime contract answer different
questions. A successful smaller machine must satisfy both, and its
blank startup must reach a configuration covered by their common
invariant. The failure just described is retained as a concrete negative
control, rather than treated as evidence for the desired state count.
